\subsection{回路格仿射扭量、粗粒化缺口与环面核素支撑的统一刚性}\label{subsec:conclusion-cyclelattice-affine-uv-coarsegraining-torsion-coding-rigidity}

沿用附录 \ref{app:fold-multiplicity} 关于回路格
\[
\Lambda(G):=\ker\partial,
\qquad
r:=\rank H_1(G;\ZZ),
\]
加权二次型 \(Q_\rho\)、Gram 矩阵 \(\Gamma_\rho\)、加权生成树多项式 \(\tau_{1/\rho}(G)\) 以及局部粗粒化 \(G\mapsto G/\pi\) 的记号。附录 I.127--I.142 已分别给出了：仿射 \(\Theta\) 函数的 Poisson 反演、固定散度能量壳的格点主项、局部粗粒化的回路维数损失、回路判别式的矩阵树恒等式与 Lee--Yang 相位分解，以及直和回路格的球截断与描述长度下界。把这些接口继续合并之后，还可抽出一条更统一的终端链：仿射扭量在紫外极限中完全退居到指数小尾项；固定散度扇区共享同一微正则熵主项；粗粒化紫外缺口由内部回路数与生成树复杂度比严格分解；而同一个回路判别式又同时控制 \(\Theta\) 前因子、Kirchhoff 复杂度、Lee--Yang 相位谱与整数环面核的素支撑。

\begin{theorem}[回路格仿射扭量的紫外相位失明]\label{thm:conclusion-cyclelattice-affine-torsor-uv-blindness}
\leanverified{paper\_conclusion\_cyclelattice\_affine\_torsor\_uv\_blindness}
令 \(G\) 为有限连通无自环多重图，\(\rho:E(G)\to(0,\infty)\) 为正权，并固定一个仿射扭量类
\[
\psi_0\in \Lambda(G)_\RR/\Lambda(G).
\]
定义对应的仿射 \(\Theta\) 函数
\[
\Theta_{G,\rho}^{(\psi_0)}(t)
:=
\sum_{\psi\in \Lambda(G)}
\exp\!\bigl(-\pi t\,Q_\rho(\psi+\psi_0)\bigr).
\]
则存在常数 \(c_{G,\rho}>0\)，使得当 \(t\downarrow 0\) 时有
\[
\Theta_{G,\rho}^{(\psi_0)}(t)
=
t^{-r/2}\,(\det\Gamma_\rho)^{-1/2}\Bigl(1+O\bigl(e^{-c_{G,\rho}/t}\bigr)\Bigr),
\]
并且主项对 \(\psi_0\) 完全独立。特别地，
\[
\frac{\Theta_{G,\rho}^{(\psi_0)}(t)}{\Theta_{G,\rho}(t)}
=
1+O\bigl(e^{-c_{G,\rho}/t}\bigr).
\]
\end{theorem}
\begin{proof}
由推论 \ref{cor:graph-cycle-affine-torsor-theta-poisson-phase}，
\[
\Theta_{G,\rho}^{(\psi_0)}(t)
=
t^{-r/2}\,(\det\Gamma_\rho)^{-1/2}
\sum_{\xi\in\Lambda(G)^\vee}
\exp\!\left(-\pi\frac{\langle \xi,\Gamma_\rho^{-1}\xi\rangle}{t}\right)
e^{2\pi i\langle \xi,\psi_0\rangle}.
\]
其中 \(\xi=0\) 项恒等于 \(1\)。其余各项的模长均由
\[
\exp\!\left(-\pi\frac{\langle \xi,\Gamma_\rho^{-1}\xi\rangle}{t}\right)
\le
\exp(-c_{G,\rho}/t),
\qquad
c_{G,\rho}:=
\pi\min_{\xi\in\Lambda(G)^\vee\setminus\{0\}}
\langle \xi,\Gamma_\rho^{-1}\xi\rangle
>0
\]
控制，于是得到第一式。另一方面，定理 \ref{thm:graph-cycle-lattice-theta-modular-inversion} 对无平移情形给出
\[
\Theta_{G,\rho}(t)
=
t^{-r/2}\,(\det\Gamma_\rho)^{-1/2}\Bigl(1+O\bigl(e^{-c'_{G,\rho}/t}\bigr)\Bigr).
\]
取 \(c_{G,\rho}:=\min(c_{G,\rho},c'_{G,\rho})\) 并比较两式即可得到比值估计。证毕。
\end{proof}

\begin{theorem}[固定散度扇区的微正则等熵律]\label{thm:conclusion-cyclelattice-fixed-divergence-equientropy}
\leanverified{paper\_conclusion\_cyclelattice\_fixed\_divergence\_equientropy}
令 \(G\) 为有限连通无自环多重图，\(w:E(G)\to(0,\infty)\) 为正权。对任意 \(\delta\in\operatorname{im}\partial\)，定义
\[
N_{\delta,w}(E)
:=
\#\{\phi\in C_1(G;\ZZ):\ \partial\phi=\delta,\ \mathcal E_w(\phi)\le E\}.
\]
则对任意 \(\delta,\delta'\in\operatorname{im}\partial\)，都有
\[
\frac{N_{\delta,w}(E)}{N_{\delta',w}(E)}
=
1+O(E^{-1/2})
\qquad(E\to\infty),
\]
并且
\[
\log N_{\delta,w}(E)
=
\frac r2\log E-\frac12\log\det A+\log\operatorname{vol}(B_r)+o(1),
\]
其中 \(A=T^{\mathsf T}\diag(w_e)T\) 为任一整数回路基上的 Gram 矩阵，故右端主项与 \(\delta\) 无关。
\end{theorem}
\begin{proof}
定理 \ref{thm:graph-energy-shell-lattice-counting} 给出
\[
N_{\delta,w}(E)
=
\frac{\operatorname{vol}(B_r)}{\sqrt{\det A}}\,E^{r/2}
+O\!\left(E^{(r-1)/2}\right),
\]
其中 \(A\) 仅依赖于回路格与权重 \(w\)，与所选扇区 \(\delta\) 无关。于是任意 \(\delta,\delta'\) 的主项常数相同，相除便得
\[
\frac{N_{\delta,w}(E)}{N_{\delta',w}(E)}
=
1+O(E^{-1/2}).
\]
对数式由同一渐近式直接推出。证毕。
\end{proof}

\begin{theorem}[局部粗粒化的紫外自由能缺口分解]\label{thm:conclusion-cyclelattice-coarsegraining-uv-gap}
\leanverified{paper\_conclusion\_cyclelattice\_coarsegraining\_uv\_gap}
令 \(G\) 为有限连通无自环多重图，\(\pi:V(G)\twoheadrightarrow Y\) 满足定理 \ref{thm:graph-coarsegraining-cycle-rank-decomposition} 的两条连通性条件。记
\[
G_y:=G[\pi^{-1}(y)],
\qquad
\Delta_\pi:=\sum_{y\in Y}\beta_1(G_y).
\]
则存在常数 \(c>0\)，使得当 \(t\downarrow 0\) 时有
\[
\log \Theta_G(t)-\log \Theta_{G/\pi}(t)
=
\frac{\Delta_\pi}{2}\log\frac1t
-\frac12\log\frac{\tau(G)}{\tau(G/\pi)}
+O(e^{-c/t}).
\]
因此，局部粗粒化的紫外自由能缺口被严格分解为一个仅由被塌缩内部回路总数控制的普适对数项，与一个仅由生成树复杂度比控制的行列式修正项。
\end{theorem}
\begin{proof}
定理 \ref{thm:graph-coarsegraining-cycle-rank-decomposition} 给出
\[
\beta_1(G/\pi)=\beta_1(G)-\Delta_\pi.
\]
另一方面，由定理 \ref{thm:hypercube-cycle-lattice-theta-asymptotic}，
\[
\log\Theta_H(t)
=
\frac{\beta_1(H)}2\log\frac1t-\frac12\log\tau(H)+O(e^{-c_H/t})
\qquad (H=G,G/\pi).
\]
两式相减，并取 \(c:=\min(c_G,c_{G/\pi})\)，即得
\[
\log \Theta_G(t)-\log \Theta_{G/\pi}(t)
=
\frac{\beta_1(G)-\beta_1(G/\pi)}2\log\frac1t
-\frac12\log\frac{\tau(G)}{\tau(G/\pi)}
+O(e^{-c/t}),
\]
再代入 \(\beta_1(G)-\beta_1(G/\pi)=\Delta_\pi\) 即可。证毕。
\end{proof}

\begin{corollary}[粗粒化恢复的 Gödel 辅助比特税]\label{cor:conclusion-cyclelattice-coarsegraining-godel-bit-tax}
在定理 \ref{thm:conclusion-cyclelattice-coarsegraining-uv-gap} 的设定下，对任意有限连通图 \(H\) 与 \(R>0\) 记
\[
\Lambda(H;R):=\{z\in\Lambda(H):\ |z|\le R\}.
\]
则当 \(R\to\infty\) 时有
\[
\log |\Lambda(G;R)|-\log |\Lambda(G/\pi;R)|
=
\Delta_\pi\log R
-\frac12\log\frac{\tau(G)}{\tau(G/\pi)}
+O(1).
\]
从而，对任意单射编码
\[
\Lambda(G;R)\hookrightarrow \Lambda(G/\pi;R)\times \{0,1\}^{b_R},
\]
都必须满足
\[
b_R
\ge
\Delta_\pi\log_2 R
-\frac12\log_2\frac{\tau(G)}{\tau(G/\pi)}
+O(1).
\]
\end{corollary}
\begin{proof}
由推论 \ref{cor:hypercube-cycle-lattice-counting-godel-log}，
\[
\log |\Lambda(H;R)|
=
\beta_1(H)\log R-\frac12\log\tau(H)+O(1)
\qquad(H=G,G/\pi).
\]
相减并代入
\(
\beta_1(G)-\beta_1(G/\pi)=\Delta_\pi
\)
即得第一式。若存在所述单射编码，则
\[
|\Lambda(G;R)|\le |\Lambda(G/\pi;R)|\,2^{b_R}.
\]
两边取二进对数，再代入前式即可得到比特下界。证毕。
\end{proof}

\begin{theorem}[回路判别式的三重统一]\label{thm:conclusion-cyclelattice-discriminant-triple-unity}
\leanverified{paper\_conclusion\_cyclelattice\_discriminant\_triple\_unity}
令 \(G\) 为有限连通无自环多重图，\(\rho:E(G)\to(0,\infty)\) 为正权，并定义回路判别式
\[
\mathscr D_\rho(G):=\det\Gamma_\rho.
\]
对任意满足 \(D\ge \lambda_{\max}(L_{1/\rho})/2\) 的常数 \(D\)，取命题 \ref{prop:graph-cycle-lattice-leyang-phase-polynomial} 的角变量 \(\theta_2,\dots,\theta_{|V|}\in(0,\pi]\)。则有精确恒等式
\[
\mathscr D_\rho(G)
=
\Bigl[\lim_{t\downarrow0} t^{r/2}\Theta_{G,\rho}(t)\Bigr]^{-2}
=
\Bigl(\prod_{e\in E(G)}\rho_e\Bigr)\tau_{1/\rho}(G)
=
\Bigl(\prod_{e\in E(G)}\rho_e\Bigr)\frac{(4D)^{|V|-1}}{|V|}
\prod_{j=2}^{|V|}\sin^2\!\frac{\theta_j}{2}.
\]
因此，回路格 \(\Theta\) 函数的紫外前因子、加权生成树多项式与 Lee--Yang 型单位圆相位谱实际上编码的是同一回路判别式。
\end{theorem}
\begin{proof}
由定理 \ref{thm:graph-cycle-lattice-theta-modular-inversion}，
\[
\Theta_{G,\rho}(t)\sim t^{-r/2}(\det\Gamma_\rho)^{-1/2},
\]
故
\[
\mathscr D_\rho(G)
=
\Bigl[\lim_{t\downarrow0} t^{r/2}\Theta_{G,\rho}(t)\Bigr]^{-2}.
\]
由定理 \ref{thm:graph-cycle-lattice-weighted-discriminant}，
\[
\mathscr D_\rho(G)=\Bigl(\prod_e\rho_e\Bigr)\tau_{1/\rho}(G).
\]
再由命题 \ref{prop:graph-cycle-lattice-leyang-phase-polynomial}，
\[
\tau_{1/\rho}(G)
=
\frac{(4D)^{|V|-1}}{|V|}\prod_{j=2}^{|V|}\sin^2\!\Bigl(\frac{\theta_j}{2}\Bigr).
\]
将后两式合并即得结论。证毕。
\end{proof}

\begin{theorem}[紫外前因子与环面核的素支撑刚性]\label{thm:conclusion-cyclelattice-uvprefactor-torus-prime-support}
\leanverified{paper\_conclusion\_cyclelattice\_uvprefactor\_torus\_prime\_support}
在定理 \ref{thm:conclusion-cyclelattice-discriminant-triple-unity} 的设定下，进一步假设 \(\rho_e=w_e\in\ZZ_{>0}\)，并取一个整数回路基使
\[
A:=T^{\mathsf T}\diag(w_e)T\in M_r(\ZZ).
\]
考虑环面自同态
\[
F_A:\TT^r\to\TT^r,
\qquad
x\mapsto Ax.
\]
则对任意素数 \(\ell\)，以下四个条件等价：
\[
\ell\mid \det A,
\qquad
\ker(F_A)[\ell^\infty]\neq 0,
\qquad
\ell\mid \Bigl(\prod_{e\in E(G)}w_e\Bigr)\tau_{1/w}(G),
\qquad
\ell\mid \Bigl[\lim_{t\downarrow0} t^{r/2}\Theta_{G,w}(t)\Bigr]^{-2}.
\]
换言之，环面核的全部素支撑既可由加权生成树复杂度恢复，也可由回路格 \(\Theta\) 函数的紫外前因子恢复；这两者与整数 Gram 矩阵 \(A\) 的 Smith 型素谱完全一致。
\end{theorem}
\begin{proof}
由推论 \ref{cor:graph-cycle-lattice-determinant-prime-support}，
\[
\ell\mid \det A
\iff
\ker(F_A)[\ell^\infty]\neq 0.
\]
另一方面，在当前所选整数回路基下，
\[
\det A=\det\Gamma_w.
\]
再由定理 \ref{thm:conclusion-cyclelattice-discriminant-triple-unity}，
\[
\det\Gamma_w
=
\Bigl(\prod_{e\in E(G)}w_e\Bigr)\tau_{1/w}(G)
=
\Bigl[\lim_{t\downarrow0} t^{r/2}\Theta_{G,w}(t)\Bigr]^{-2}.
\]
故四个条件完全等价。证毕。
\end{proof}

\begin{corollary}[独立 Stokes 扇区的张量化体积律与线性编码不可压缩性]\label{cor:conclusion-cyclelattice-directsum-coding-incompressibility}
\leanverified{paper\_conclusion\_cyclelattice\_directsum\_coding\_incompressibility}
令 \(G\) 为有限连通多重图，并记
\[
r:=\beta_1(G),
\qquad
\tau(G):=\text{生成树个数}.
\]
对任意 \(m\ge 1\)，考虑直和格
\[
\Lambda(G)^{\oplus m}
=
\underbrace{\Lambda(G)\oplus\cdots\oplus\Lambda(G)}_{m\ \text{份}}.
\]
则其球截断满足
\[
\log |\Lambda(G)^{\oplus m}(R)|
=
mr\log R-\frac m2\log\tau(G)+O(1)
\qquad(R\to\infty).
\]
因此，任意把 \(\Lambda(G)^{\oplus m}(R)\) 无损编码为二进制字符串的方案，其最小码长 \(L_m(R)\) 必满足
\[
L_m(R)\ge mr\log_2 R-\frac m2\log_2\tau(G)+O(1).
\]
特别地，独立回路扇区的可分辨体积与最小码长都严格呈线性可加，不存在能够把 \(m\) 份独立回路扇区压缩到低于线性主阶的统一编码。
\end{corollary}
\begin{proof}
推论 \ref{cor:hypercube-cycle-lattice-directsum-counting} 已给出
\[
\log |\Lambda(G)^{\oplus m}(R)|
=
mr\log R-\frac m2\log\tau(G)+O(1).
\]
再由推论 \ref{cor:hypercube-cycle-lattice-directsum-mdl}，任意有限字母表上的单射编码长度都至少等于该体积对数除以字母表对数。取二进字母表即得
\[
L_m(R)\ge mr\log_2 R-\frac m2\log_2\tau(G)+O(1).
\]
证毕。
\end{proof}

\endinput
